\documentclass[12pt,a4paper]{article}

\usepackage[utf8]{inputenc}
\usepackage[T1]{fontenc}
\usepackage[margin=2.5cm]{geometry}
\usepackage{hyperref}
\usepackage{microtype}
\usepackage{parskip}
\usepackage{lmodern}
\usepackage{amsmath}
\usepackage{amssymb}

\hypersetup{colorlinks=true, urlcolor=blue!60!black}

\pagestyle{plain}

\begin{document}

\begin{center}
  {\large\textbf{Research Proposal}} \\[6pt]
  {\large Partition Geometry as a Predictive Framework for Soft Matter Rheology
  and Differential Dynamic Microscopy} \\[6pt]
  Kundai Farai Sachikonye --- Application for predoctoral position, SoMeX Lab
  (Job-ID 5570)
\end{center}

\vspace{1em}

\subsection*{Central Question}

Can the rheological properties of yield stress fluids and the intermediate
scattering functions measured by Differential Dynamic Microscopy (DDM) be
predicted from first principles --- without empirical fitting --- using a single
unified framework based on bounded phase space geometry?

\subsection*{Background and Motivation}

Yield stress fluids (concentrated emulsions, colloidal gels, polymer networks,
biological hydrogels) are among the most practically important and theoretically
challenging soft matter systems. Their defining property --- the existence of a
critical stress $\sigma_y$ below which the material behaves as a solid and above
which it flows --- is typically treated phenomenologically: the Herschel-Bulkley
or Bingham models fit experimental flow curves but provide no prediction of $\sigma_y$
from molecular or colloidal parameters.

DDM offers a powerful experimental route to the dynamic properties of such systems:
the image structure function $D(\mathbf{q}, \Delta t)$ gives access to the
intermediate scattering function $f(\mathbf{q}, \Delta t)$ without the
sample constraints of DLS, enabling in-situ characterisation of structure and
dynamics during flow.
Yet the interpretation of DDM data in yield stress fluids remains challenging:
the spatial and temporal heterogeneity of these materials, and the optical
complexity of concentrated colloidal dispersions, make connecting $D(\mathbf{q}, \Delta t)$
to mechanical properties non-trivial.

\subsection*{Proposed Framework}

I propose to develop and test a first-principles connection between the
rheological behaviour and DDM-measured dynamics of model yield stress systems,
using the bounded phase space (BPS) framework I have developed.

\textbf{Step 1 --- Rheological predictions from partition geometry.}
In the BPS framework, every soft matter system is a bounded oscillatory system
characterised by S-entropy coordinates $(S_k, S_t, S_e) \in [0,1]^3$.
Viscosity is $\mu = \tau_c \times g$, where $\tau_c$ is the partition lag
(timescale of spectral mode transitions) and $g$ is the coupling strength
(off-diagonal elements of the spectral coherence matrix).
The yield stress $\sigma_y$ corresponds to the partition floor: the minimum
S-entropy coordinate $S_e^{\min} = 1/m > 0$ (Floor Theorem, where $m$ is the
number of spectral modes), below which the ordered phase cannot sustain coherence.
This gives a structural prediction: $\sigma_y$ is proportional to the spectral
coherence floor of the colloidal assembly, which can be computed from particle
size distribution and volume fraction before any mechanical measurement.

\textbf{Step 2 --- DDM structure function from optical partition geometry.}
The image structure function $D(\mathbf{q}, \Delta t) = 2[S(\mathbf{q}) -
\mathrm{Re}\{S(\mathbf{q}, \Delta t)\}]$, where $S(\mathbf{q}, \Delta t)$
is the dynamic structure factor.
In the BPS framework, the optical signal in a DDM experiment is governed by
the Resonant Stack theorem: the transmittance $T = 1/e$ is the universal
fixed point of the Beer-Lambert-Kirchhoff equivalence, and the image contrast
encoding particle dynamics is a projection of the partition structure function
onto the imaging wave vector $\mathbf{q}$.
This gives a predictive model for $D(\mathbf{q}, \Delta t)$ from the same
partition geometry that predicts $\mu$ and $\sigma_y$ --- a single framework
connecting rheology to DDM without separate fitting.

\textbf{Step 3 --- Experimental validation in the SoMeX Lab.}
The proposal would be tested against measurements on two model systems:
(i) Carbopol microgel suspensions (archetypal soft yield stress fluid;
well-characterised by the lab's Anton Paar MCR 702e and ShearView rheometer);
(ii) hard-sphere colloidal suspensions near the glass transition
(DDM-accessible via the lab's sCMOS camera systems and fastDDM code).
For each system, the BPS framework predicts $\sigma_y$ and the shape of
$D(\mathbf{q}, \Delta t)$ from the particle parameters alone.
Agreement or disagreement with measurement is a direct test of whether the
partition geometry captures the relevant physics.

\subsection*{Expected Outcomes and Timeline}

Year 1: Implement BPS rheology predictions for Carbopol and hard-sphere systems;
compare to ShearView measurements; extend fastDDM analysis to extract
structure functions for BPS validation.

Year 2--3: Extend to biologically relevant yield stress systems (mucus, cytoskeletal
networks, bacterial biofilm matrices) and to active matter DDM (bacterial suspensions),
where the $\oint \mathbf{J}\cdot d\mathbf{A} = 0$ conservation law for collective
motion provides an additional first-principles constraint on $f(\mathbf{q}, \Delta t)$.

The outcome, if the predictions hold, would be a validated first-principles
connection between particle-level structure and macroscopic rheological
and optical observables in soft matter systems --- a contribution both to
fundamental soft matter physics and to the practical characterisation of
biological fluids and active matter.

\end{document}
